\section*{Open questions (5 truly-open, with concrete next steps)}

\begin{enumerate}
  \item \textbf{Non-Pauli / non-unitary noise models.}
  Does the RC error-suppression exponent (paper $\sim$2.73, this work exact $\sim$2.26) hold for non-Pauli coherent noise --- e.g., amplitude-damping Kraus pairs or correlated $ZZ$ crosstalk between neighboring qubits?
  \emph{Basis:} Theorem~1 strictly requires all-Hermitian Kraus operators; single-qubit $R_Z(\varepsilon)$ is the easiest case. Non-unitary $T_1$-like noise is only partially symmetrized by 1-qubit Pauli twirl.
  \emph{Next steps:} Extend the exact twirled-channel calculation with a per-cycle noise model mixing $R_Z(\varepsilon)$ and an amplitude-damping Kraus pair of strength $\gamma$. Fit RC exponent vs $(\varepsilon,\gamma)$. Also test a 2-qubit $ZZ(\alpha)$ crosstalk twirled by the 2-qubit Pauli group. Predict: RC exponent falls from $\sim$2.3 toward 1 as the non-Hermitian fraction grows.

  \item \textbf{RC composed with dynamical decoupling (DD).}
  Is the combined coherent-noise suppression of RC and DD additive, multiplicative, or do they interfere?
  \emph{Basis:} DD suppresses low-frequency coherent noise via pulse echoes; RC suppresses per-cycle coherent noise via logical Pauli twirls. Both target coherent errors at different timescales; the paper does not combine them.
  \emph{Next steps:} Insert XY4 or CPMG DD blocks inside each iterative-PE cycle of the qiskit-aer harness. Repeat the exact-superoperator sweep as a function of $(\varepsilon, T_{\text{DD}}, N_r)$. Metric: RC-only slope vs DD-only slope vs combined slope on the same coherent-noise sweep. Free endpoint: local qiskit-aer.

  \item \textbf{Chemistry ground-state phase estimation.}
  Does the RC exponent $\sim$2.7 hold when the estimated phase is a chemical ground-state eigenphase (H$_2$, LiH) rather than a synthetic Floquet quasi-energy?
  \emph{Basis:} Paper Fig.~3 uses a synthetic 10-qubit Floquet unitary and Fig.~4 uses Shor. Chemistry phases have near-degenerate eigenvalue clusters that may interact with RC's stochastic residual differently.
  \emph{Next steps:} Build the H$_2$/STO-3G Bravyi-Kitaev qubit Hamiltonian in \texttt{qiskit-nature}, Trotterize $e^{-iHt}$, run bare iterative-PE and RC on the ground-state overlap. Fit the error exponent under coherent Trotter noise. Free endpoint: local \texttt{qiskit-aer}.

  \item \textbf{Hardware-runtime overhead of per-shot Pauli twirling.}
  On real IBM/IonQ hardware, once compiler-side twirl injection and per-shot recompilation are counted, does the RC advantage survive at fixed wall-clock budget?
  \emph{Basis:} This replication is simulator-only. RC requires $N_r$ distinct compiled circuits per depth $L$; each recompile hits the transpiler and pulse-schedule build. At small total-shot budgets, compile time may dominate.
  \emph{Next steps:} Run iterative-PE + RC on a free IBM Q Experience open-plan backend (e.g.\ \texttt{ibm\_brisbane}) with $N_r\in\{1,10,100\}$ at fixed total shot budget $T$. Fair comparison: bare-error at $T$ shots vs RC-error at $T/N_r$ shots. Predict: RC advantage shrinks at small $T$.

  \item \textbf{Adaptive shot allocation across RC realizations.}
  Is there a Cramer-Rao-optimal $(N_s(L), N_r(L))$ allocation that outperforms uniform-per-depth by concentrating twirl realizations where the RC bias residual is largest?
  \emph{Basis:} RC residual scales as $1/\sqrt{N_s N_r} + \varepsilon^k$-bias. At small $L$ bias is negligible; at large $L$ aliasing of $\cos(2\varphi L)$ makes small residuals dominate. Uniform allocation is provably suboptimal.
  \emph{Next steps:} Formulate a variance cost $C(L, N_s, N_r; \varepsilon)$ from the Aer-sampled likelihood; minimize under a total budget $\sum_L N_s(L) N_r(L) = T$. Compare a Cramer-Rao-optimal allocation against uniform on the exact-channel sweep. Metric: $\varphi$-error at fixed $T$.
\end{enumerate}
